%% main.tex — Citation Manipulation Study
%% Target: Quantitative Science Studies (QSS), MIT Press
%% Compile: pdflatex main && bibtex main && pdflatex main && pdflatex main
%% ------------------------------------------------------------------
\documentclass[12pt,a4paper]{article}

% ── Packages ────────────────────────────────────────────────────────
\usepackage[T1]{fontenc}
\usepackage[utf8]{inputenc}
\usepackage{microtype}                 % microtypography
\usepackage{mathptmx}                  % Times-compatible serif
\usepackage[margin=2.5cm]{geometry}
\usepackage{setspace}
\doublespacing

% Figures
\usepackage{graphicx}
\usepackage{float}
\usepackage[labelfont=bf, labelsep=period,
            font=small, justification=raggedright,
            singlelinecheck=false]{caption}
\usepackage{subcaption}
\usepackage{booktabs}                  % professional tables
\usepackage{amsmath}
\usepackage{amssymb}                   % \blacklozenge etc.
\usepackage{xcolor}
\definecolor{caseRed}{HTML}{b2182b}
\definecolor{ctrlBlue}{HTML}{2166ac}

% Hyperlinks
\usepackage[colorlinks=true,
            linkcolor=black,
            citecolor=black,
            urlcolor=blue]{hyperref}

% References
\usepackage[round,sort,comma,authoryear]{natbib}
\bibliographystyle{abbrvnat}

% Figure path
\graphicspath{{analysis_results_v4/figs/}}

% ── Metadata ────────────────────────────────────────────────────────
\title{\textbf{Co-authorship Citation Syndicates:\\
       Network Evidence of Coordinated Citation Inflation}}

\author{
  \textit{[Author names omitted for blind review]}
}

\date{}

%% =================================================================
\begin{document}
\maketitle

% ── Abstract ────────────────────────────────────────────────────────
\begin{abstract}
\noindent
We investigate systematic co-author citation inflation in academic
publishing using a matched-pair design across 9{,}431 author dyads
spanning five subject fields.
Applying Isolation Forest anomaly detection to a 13-dimensional
graph-topology feature space, we identify 189 outlier authors
($\approx1\%$ contamination rate) whose co-author citation patterns
deviate substantially from matched controls.
Outlier authors exhibit a co-author citation rate $6.7\times$ higher
than controls (0.0040 vs.\ 0.0006, $p < 0.0001$, permutation test)
and form a densely connected syndicate of 23 members with an
internal citation density of 0.29.
Tier assortativity ($r = 0.707$) confirms strong homophily in
citation behaviour.
Linear discriminant analysis achieves clean bimodal separation,
with Outgoing HHI (29\%) and Journal Endogamy (27\%) as the leading
discriminant features.
These results provide quantitative evidence of coordinated citation
inflation and propose a reproducible detection framework applicable
to large-scale bibliometric surveillance.

\medskip\noindent\textbf{Keywords:}
citation manipulation; co-author citation rings; anomaly detection;
network assortativity; bibliometrics; research integrity
\end{abstract}

\newpage

%% =================================================================
\section{Introduction}\label{sec:intro}

The integrity of citation metrics underpins hiring, promotion, and
funding decisions in contemporary academia.
Citation counts and derived indicators such as the $h$-index are
increasingly gamed through self-citation inflation, gift authorship,
and \emph{citation cartels}---coordinated groups of authors who
systematically inflate one another's citation counts
\citep[e.g.,][]{vannoorden2013, fister2016}.
While isolated cases have been documented anecdotally, scalable
automated detection of such syndicates remains methodologically
underdeveloped.

This paper addresses two questions:
\begin{enumerate}
  \item Can graph-topology features derived from a large-scale
        citation network discriminate between ``inflating'' and
        ``normal'' authors matched on observable academic profile?
  \item Do detected outliers cluster into identifiable citation
        syndicates?
\end{enumerate}

We answer both questions affirmatively using a combination of
matched-pair design, Isolation Forest outlier detection, Random
Forest classification, and linear discriminant analysis applied to a
population of 18{,}862 authors across five broad subject fields.

%% =================================================================
\section{Data and Methods}\label{sec:methods}

\subsection{Data Sources}

Author profiles, publications, and citation linkages are drawn from
two SQLite databases.
\texttt{impact.db} contains raw tables (\texttt{works},
\texttt{work\_authors}, \texttt{work\_references}) produced by
processing OpenAlex metadata.
\texttt{rolap.db} contains derived analysis tables including
pre-computed behaviour metrics (\texttt{author\_behavior\_metrics}),
anomaly signals (\texttt{citation\_anomalies}), and the
weighted directed citation network
(\texttt{citation\_network\_final}).

\subsection{Matched-Pair Design}

For each subject field $s$ we stratify authors by journal-impact
percentile and construct 1:1 matched pairs of Bottom-Tier
(\emph{Case}) and Top-Tier (\emph{Control}) authors.
Matching is performed on career age, publication count, and subject
field, yielding 9{,}431 pairs across five subject fields
(Medicine, Engineering, Natural Sciences, Social Sciences,
Arts~\& Humanities).

\subsection{Feature Engineering}

Thirteen graph-topology features are computed from the weighted
directed citation network:
co-author citation rate, self-citation rate, local clustering
coefficient, normalised triangle count, citation balance,
reciprocity rate, outgoing HHI, clique strength, PageRank,
k-core number, incoming citation entropy, incoming HHI, and journal
endogamy rate.
All features are scaled with a robust scaler prior to modelling.

\subsection{Outlier Detection}

Isolation Forest \citep{liu2008} is applied with contamination rate
$\gamma = 0.01$, yielding 189 outlier authors.
We verify detection robustness using permutation testing
($B = 10{,}000$) on the mean paired co-author citation rate difference.

\subsection{Discriminant Analysis}

A Random Forest classifier (300 trees, max depth 8,
out-of-bag evaluation) is trained on the 13-feature matrix to rank
feature importances.
A one-dimensional Linear Discriminant Analysis (LDA) projection is
used to visualise tier separability.

%% =================================================================
\section{Results}\label{sec:results}

\subsection{Co-author Citation Gap Across Fields}

Figure~\ref{fig:forest} shows the mean paired difference in
co-author citation rate (Case $-$ Control) for each subject field
and for the overall sample.
All fields exhibit a positive gap, confirming that Bottom-Tier
authors cite co-authors at higher rates than matched Top-Tier peers.
The overall effect is $\Delta = 0.0033$ (95\% CI includes zero only
for Arts~\&~Humanities), with Engineering showing the largest
field-specific gap ($\Delta = 0.0048$).

\begin{figure}[H]
  \centering
  \includegraphics[width=\textwidth]{Figure1_ForestPlot}
  \caption{%
    \textbf{(a) Co-author citation gap by subject field.}
    Forest plot of mean paired difference
    $\Delta = \bar{r}_{\text{Case}} - \bar{r}_{\text{Control}}$
    in co-author citation rate, with 95\% confidence intervals.
    Marker $\bullet$ = field estimate; $\blacklozenge$ = overall.
    $\Delta$ labels report rounded effect sizes.
    The dashed vertical line marks zero (no difference);
    the pale red shading highlights the Case-higher region.
  }
  \label{fig:forest}
\end{figure}

\subsection{Outlier Behavioural Fingerprint}

Figure~\ref{fig:radar} characterises the 189 outliers relative to
normal authors across six behavioural dimensions expressed as
fold-change ratios.
Co-author citation rate shows the most extreme elevation
($94.3\times$ above the normal mean), followed by reciprocity
($68.4\times$) and normalised triangles ($43.4\times$).
This profile is consistent with a strategy of dense, bilateral
co-author citation exchange within a tightly connected cluster.

\begin{figure}[H]
  \centering
  \includegraphics[width=0.6\textwidth]{Figure2a_Fingerprint}
  \caption{%
    \textbf{(b) Outlier behavioural fingerprint.}
    Radar chart of outlier fold-change ratios relative to normal
    authors across six citation-behaviour metrics.
    Values on each spoke indicate the outlier-to-normal mean ratio.
    The dashed ring marks the baseline (1$\times$ normal).
  }
  \label{fig:radar}
\end{figure}

\subsection{Outlier Citation Syndicate}

Figure~\ref{fig:network} visualises the largest connected component
among mutually citing outlier authors.
The 23-member group exhibits an internal citation density of 0.29
--- far above the network-wide baseline --- with one hub author
(salmon node) mediating $\approx$29\% of intra-group traffic.
Edge widths are proportional to citation weight;
the prevalence of thick bidirectional edges is consistent with
coordinated mutual citation inflation rather than organic academic
exchange.

\begin{figure}[H]
  \centering
  \includegraphics[width=\textwidth]{Figure4_Network}
  \caption{%
    \textbf{(c) Largest outlier citation syndicate ($n = 23$).}
    Force-directed layout of the internal citation network.
    Node size is proportional to betweenness centrality;
    \colorbox{caseRed!20}{\strut salmon} = hub (maximum betweenness);
    {\color{caseRed}red} = syndicate member.
    Edge width scales with aggregate citation weight.
    The inset box summarises graph statistics.
  }
  \label{fig:network}
\end{figure}

\subsection{Discriminant Feature Importances}

Figure~\ref{fig:importance} reports Random Forest feature importances
for the Case--Control classification task.
Outgoing HHI (29\%) and Journal Endogamy (27\%) are the two
dominant features, indicating that concentration of outgoing
citations towards a small set of recipients, and preferential
citation of journals where the author also publishes,
are the strongest discriminators.
Higher-importance features are shown in {\color{caseRed}red};
lower-importance features in {\color{ctrlBlue}blue}.

\begin{figure}[H]
  \centering
  \includegraphics[width=\textwidth]{Figure8_Feature_Importance}
  \caption{%
    \textbf{(d) Random Forest feature importances.}
    Horizontal bars show mean decrease in impurity for each of the
    13 graph-topology features, sorted ascending.
    {\color{caseRed}Red} bars denote features at or above the median
    importance; {\color{ctrlBlue}blue} bars denote below-median.
    Inline labels give raw importance score and share of total.
  }
  \label{fig:importance}
\end{figure}

\subsection{Tier Separability}

Figure~\ref{fig:lda} shows the 1-D LDA projection of all 18{,}862
authors onto the single discriminant axis.
The two tier distributions are clearly offset, with Case authors
skewed towards negative scores (cohesion-driven structure) and
Control authors towards positive scores (authority-driven structure).
The bimodal shape of the Case distribution indicates that the sample
contains two sub-populations: a large group with moderate citation
patterns, and a smaller, more extreme group corresponding to
flagged outliers.

\begin{figure}[H]
  \centering
  \includegraphics[width=0.75\textwidth]{Figure9_LDA_Separation}
  \caption{%
    \textbf{(e) Tier separability: LDA projection.}
    Kernel density estimates of the 1-D linear discriminant scores
    for {\color{caseRed}Bottom-Tier (Case)} and
    {\color{ctrlBlue}Top-Tier (Control)} authors.
    Dotted vertical lines mark within-tier medians.
    Directional labels indicate the dominant behavioural regime
    at each tail of the discriminant axis.
  }
  \label{fig:lda}
\end{figure}

\subsection{Citation Tier Homophily}

Figure~\ref{fig:mixing} displays the row-normalised citation mixing
matrix.
Top-Tier authors direct 94.7\% of their citations to fellow
Top-Tier authors, while Bottom-Tier authors direct 73.4\% of
citations within their own tier.
The resulting tier assortativity coefficient $r = 0.707$ indicates
strong homophily: authors overwhelmingly cite within their own
prestige stratum.
The asymmetry between rows reveals that Top-Tier authors are
substantially more insular, whereas Bottom-Tier authors show
greater cross-tier citation (26.6\% to Top-Tier), consistent with
aspirational or strategic upward citing.

\begin{figure}[H]
  \centering
  \includegraphics[width=0.72\textwidth]{Figure10_Mixing_Matrix}
  \caption{%
    \textbf{(f) Citation mixing matrix ($r = 0.707$, strong homophily).}
    Row-normalised proportion of citations directed within and
    across tiers.
    Cell values show the conditional probability that an author
    in the row tier cites an author in the column tier.
    The Brewer RdBu colormap is centred at 0.5.
    Below: diagonal average (84.0\%) summarises overall within-tier
    preference.
  }
  \label{fig:mixing}
\end{figure}

%% =================================================================
\section{Discussion}\label{sec:discussion}

The results provide converging evidence from four complementary
analytical lenses --- paired hypothesis testing, anomaly detection,
network analysis, and machine learning --- for the existence of
coordinated citation inflation in the studied population.

Several features of the outlier profile are theoretically
informative.
The extreme elevation in \emph{outgoing} HHI ($\approx$2.2$\times$
normal) indicates not merely higher citation rates but highly
concentrated targeting of the same small set of co-authors ---
a signature of strategic rather than organic citation behaviour.
The reciprocity elevation ($68.4\times$) confirms that these
exchanges are largely bilateral: if outlier A cites
outlier B, B almost certainly cites A in return.
Combined with the syndicate network structure (density = 0.29,
single dominant hub), this pattern fits the theoretical model of a
\emph{citation cartel} \citep{fister2016}.

Journal endogamy (27\% of feature importance) is a novel signal
not previously highlighted in the citation-manipulation literature:
outlier authors preferentially cite papers published in the same
journals they publish in, suggesting coordinated within-journal
mutual citation as a secondary inflation mechanism.

The strong tier homophily ($r = 0.707$) has implications for the
broader citation network: if bottom-tier authors largely cite each
other, their impact metrics are decoupled from genuine academic
influence and instead reflect an internal credit system.
This finding underscores the value of network-level assortativity
as a diagnostic complementing individual author-level metrics.

\subsection*{Limitations}

Our outlier labels are generated by Isolation Forest with a fixed
contamination prior; the 1\% rate is conservative and may exclude
some moderate inflators.
The dataset covers only authors with ORCID linkages and
OpenAlex coverage, potentially under-representing fields with lower
metadata completeness.
Causal claims about intentionality cannot be established from
observational data alone.

%% =================================================================
\section{Conclusion}\label{sec:conclusion}

We demonstrate that a 13-dimensional graph-topology feature space,
derived entirely from a large-scale weighted citation network,
supports accurate discrimination of citation-inflating authors from
matched controls.
The detected outliers form a structurally coherent syndicate
characterised by dense reciprocal citation, extreme concentration,
and strong assortativity.
The pipeline is fully automated and scalable, making it suitable
for deployment as part of ongoing bibliometric surveillance by
journal editors, funding bodies, or research integrity offices.

%% =================================================================
\section*{Data Availability}

All figures are generated by the open-source pipeline
\texttt{diagrams.py} and reproducible from
\texttt{rolap.db} / \texttt{impact.db}.
Code is available at \url{https://github.com/[anonymised]}.

%% =================================================================
\bibliographystyle{abbrvnat}

% ── Inline mini-bibliography (no external .bib required) ────────────
\begin{thebibliography}{9}

\bibitem[Fister et~al.(2016)]{fister2016}
Fister,~I., Fister,~I.~Jr., \& Perc,~M. (2016).
Toward the discovery of citation cartels in citation networks.
\textit{Frontiers in Physics}, 4, 49.

\bibitem[Liu et~al.(2008)]{liu2008}
Liu,~F.~T., Ting,~K.~M., \& Zhou,~Z.-H. (2008).
Isolation forest.
In \textit{Proceedings of the 8th IEEE International Conference on
Data Mining} (pp. 413--422).

\bibitem[Van~Noorden(2013)]{vannoorden2013}
Van~Noorden,~R. (2013).
Brazilian citation scheme outed.
\textit{Nature}, 500(7464), 510--511.

\end{thebibliography}

\end{document}
